\documentclass[11pt]{article}
\usepackage[margin=1in]{geometry}
\usepackage{amsmath,amssymb,amsthm}
\usepackage{hyperref}
\usepackage{enumitem}

\newtheorem{theorem}{Theorem}
\newtheorem{lemma}[theorem]{Lemma}
\newtheorem{corollary}[theorem]{Corollary}
\theoremstyle{definition}
\newtheorem{definition}[theorem]{Definition}
\theoremstyle{remark}
\newtheorem{remark}[theorem]{Remark}

\usepackage{xurl}
\let\oldus\_
\renewcommand{\_}{\oldus\allowbreak}

\title{A disproof of Conjecture 00000007785\\[2pt]
\large The cube has Mahler product $4^n/n! < 4^n$}
\date{}

\begin{document}
\sloppy
\maketitle

\section{The conjecture}

Conjecture 00000007785 reads (English version):

\begin{quote}
Definition: The Mahler product $M(K) = \mathrm{vol}(K)\,\mathrm{vol}(K^\circ)$ with polar $K^\circ$;
the Mahler conjecture asserts $M(K) \ge 4^n$ in the symmetric case. Conjecture: The symmetric
$n$-dimensional extremal bodies are the Hanner polytopes (equality cases), with strict inequality
outside finite products of the Hanner family and finitely many low-dimensional exceptions
($n\le 3$); the nonsymmetric case $M(K) \ge (n+1)^{n+1}/(n!)^2$ is attained uniquely by the
simplex; and both statements have stability versions: $M(K) \le 4^n(1+\varepsilon)$ implies $K$ is
$O(\varepsilon^{1/2})$ close to a Hanner body.
\end{quote}

The Chinese version states the same formulas: the symmetric bound $M(K)\ge 4^n$, the Hanner
polytopes as the equality cases, and the stability hypothesis $M(K)\le 4^n(1+\varepsilon)$.

\paragraph{Reading.} We read the text literally. Its symmetric clause contains two claims about
origin-symmetric convex bodies $K\subset\mathbb{R}^n$:
\begin{enumerate}[label=(\alph*)]
\item the inequality $M(K)\ge 4^n$;
\item the Hanner polytopes are the equality cases of this bound, so every Hanner polytope satisfies
$M(K)=4^n$.
\end{enumerate}
We refute both (a) and (b) in every dimension $n\ge 2$. The conjecture is a conjunction, so it is
false. We do not address the nonsymmetric clause $(n+1)^{n+1}/(n!)^2$ or the stability clause,
and we make no claim about the statement obtained by replacing $4^n$ with $4^n/n!$ (which is the
classical Mahler conjecture; see Remark~\ref{rem:const}).

\section{Definitions and conventions}

We work in $\mathbb{R}^n$, modelled in Lean as \texttt{Fin n -> R}, with the standard dot product
$x\cdot y=\sum_i x_iy_i$ and Lebesgue measure \texttt{volume} (the product measure, which is
Euclidean volume). Volumes take values in $[0,\infty]$.

\begin{definition}
A \emph{convex body} is a compact convex set with non-empty interior. Wikipedia (Mahler volume):
``A convex body in Euclidean space is defined as a compact convex set with non-empty interior.''
It is \emph{origin-symmetric} if $-x\in K$ whenever $x\in K$. (Lean:
\texttt{IsSymmetricConvexBody}.)
\end{definition}

\begin{definition}
The polar of $K$ is $K^\circ=\{y : x\cdot y\le 1 \text{ for all } x\in K\}$, and the Mahler product
is $M(K)=\mathrm{vol}(K)\,\mathrm{vol}(K^\circ)$. (Lean: \texttt{polar}, \texttt{mahler}.) This is
the definition in the Wikipedia article ``Mahler volume''.
\end{definition}

\begin{definition}
The Hanner polytopes are generated recursively. Wikipedia (Hanner polytope) gives the rules
``A line segment is a one-dimensional Hanner polytope.'', ``The Cartesian product of every two
Hanner polytopes is another Hanner polytope, whose dimension is the sum of the dimensions of the
two given polytopes.'' and ``The dual of a Hanner polytope is another Hanner polytope of the same
dimension.'' The Lean predicate \texttt{IsHanner n K} is the inductive closure of the segment
$[-1,1]\subset\mathbb{R}^1$ under Cartesian products (coordinates concatenated by
\texttt{Fin.castAdd}/\texttt{Fin.natAdd}) and polar duals. It uses only the symmetric segment
$[-1,1]$, so it describes a subclass of the Hanner polytopes. Showing that a body lies in this
subclass therefore shows that it is a Hanner polytope.
\end{definition}

\noindent The cube is $C_n=[-1,1]^n=\{x:|x_i|\le 1\ \forall i\}$ and the cross-polytope is
$Q_n=\{y:\sum_i|y_i|\le1\}$.

\section{The proof}

\begin{lemma}[\texttt{polar\_cube}]
$C_n^\circ = Q_n$.
\end{lemma}
\begin{proof}
If $y\in Q_n$ and $x\in C_n$, then $x\cdot y\le\sum_i|x_i||y_i|\le\sum_i|y_i|\le1$. Conversely, if
$y\in C_n^\circ$, test with $x_i=1$ when $y_i\ge0$ and $x_i=-1$ otherwise. Then $x\in C_n$ and
$x\cdot y=\sum_i|y_i|\le1$.
\end{proof}

\begin{lemma}[\texttt{volume\_cube}, \texttt{volume\_crossPolytope}]
$\mathrm{vol}(C_n)=2^n$, and for $n\ge1$, $\mathrm{vol}(Q_n)=2^n/n!$.
\end{lemma}
\begin{proof}
$C_n$ is the box $\prod_i[-1,1]$, so its volume is $2^n$ (Mathlib \texttt{Real.volume\_Icc\_pi}).
For $Q_n$, Mathlib's \texttt{volume\_sum\_rpow\_le} with $p=1$ and $r=1$ gives
$\mathrm{vol}\{x:\sum_i|x_i|\le1\} = (2\Gamma(2))^n/\Gamma(n+1)=2^n/n!$, using $\Gamma(2)=1$ and
$\Gamma(n+1)=n!$.
\end{proof}

\begin{lemma}[\texttt{mahler\_cube}, \texttt{mahler\_cube\_lt}]
For $n\ge1$, $M(C_n)=4^n/n!$. For $n\ge2$, $M(C_n)<4^n$.
\end{lemma}
\begin{proof}
$M(C_n)=2^n\cdot 2^n/n!=4^n/n!$, and $n!>1$ when $n\ge2$.
\end{proof}

\begin{lemma}[\texttt{cube\_isSymmetricConvexBody}, \texttt{cube\_isHanner}]
$C_n$ is an origin-symmetric convex body. For $n\ge1$ it is a Hanner polytope.
\end{lemma}
\begin{proof}
$C_n$ is a closed box, so it is compact and convex. It is symmetric, and it contains the open
sup-norm unit ball, so $0$ is an interior point. For the Hanner property, $C_1=[-1,1]$ is the
segment, and $C_{m+1}$ is the Cartesian product of $C_m$ and $C_1$
(\texttt{Fin.forall\_fin\_add}). Induction finishes the proof.
\end{proof}

\begin{theorem}[\texttt{C7785.conjecture\_7785\_symmetric\_clause\_false}]
For every $n\ge2$: the cube $C_n$ is a Hanner polytope and an origin-symmetric convex body with
$M(C_n)=4^n/n!<4^n$. Hence
\begin{enumerate}[label=(\roman*)]
\item it is false that $M(K)\ge4^n$ for every origin-symmetric convex body $K\subset\mathbb{R}^n$;
\item it is false that every Hanner polytope in $\mathbb{R}^n$ satisfies $M(K)=4^n$.
\end{enumerate}
\end{theorem}

Claim (ii) shows that the Hanner polytopes are not the equality cases of the bound in (a), so claim (b) fails.
Both failures occur in every dimension $n\ge2$, so the clause ``finitely many low-dimensional
exceptions ($n\le3$)'' cannot rescue the statement. That clause modifies only the strict-inequality
part, and in any case the failure happens for all $n\ge 2$. In dimension $1$ the bound reads
$M([-1,1])=4=4^1$, so there is no failure there.

\begin{remark}\label{rem:const}
The Wikipedia article ``Mahler volume'' states: ``In particular, the polar body of a cube or
hypercube is an octahedron or cross polytope.'' It gives the Mahler volume of the hypercube as
$4^n/\Gamma(n+1)$, which agrees with the lemma above. It also states: ``The shapes with the minimum
known Mahler volume are hypercubes, cross polytopes, and more generally the Hanner polytopes which
include these two types of shapes, as well as their affine transformations.'' The classical
conjecture therefore has the constant $4^n/n!$, not the $4^n$ written in the statement. The
conjecture under review writes $4^n$ in both languages, both in the bound and in the stability
hypothesis $M(K)\le4^n(1+\varepsilon)$. We refute that literal statement and claim nothing about the
corrected one, which remains open for $n\ge4$.
\end{remark}

\section{Lean formalization and axiom audit}

File \texttt{lean/Conjecture7785/Basic.lean} (namespace \texttt{C7785}, 156 lines, only
\texttt{import Mathlib}, Lean/Mathlib v4.33.1). Main theorem:
\begin{verbatim}
theorem conjecture_7785_symmetric_clause_false (n : N) (hn : 2 <= n) :
    (IsHanner n (cube n) /\ IsSymmetricConvexBody (cube n) /\
        mahler (cube n) = ENNReal.ofReal (4 ^ n / n !) /\ mahler (cube n) < 4 ^ n) /\
      ~ (forall K, IsSymmetricConvexBody K -> (4 : ENNReal) ^ n <= mahler K) /\
      ~ (forall K, IsHanner n K -> mahler K = 4 ^ n)
\end{verbatim}
\texttt{lake build} succeeds, and \texttt{\#print axioms
C7785.conjecture\_7785\_symmetric\_clause\_false} reports only
\texttt{[propext, Classical.choice, Quot.sound]}. The file contains no \texttt{sorry} and no
\texttt{native\_decide}.

\end{document}
